\documentclass[14pt,a4paper]{extarticle}

\usepackage[T2A]{fontenc}
\usepackage[utf8]{inputenc}
\usepackage[russian]{babel}
\usepackage[a4paper,margin=2cm]{geometry}
\usepackage{amsmath,amssymb}
\usepackage{graphicx}
\usepackage{booktabs}
\usepackage{longtable}
\usepackage{float}
\usepackage{array}
\usepackage{caption}
\usepackage{hyperref}
\usepackage{setspace}

\onehalfspacing
\graphicspath{{../outputs/lab3/plots/}}

\hypersetup{
    colorlinks=true,
    linkcolor=black,
    urlcolor=blue
}

\newcommand{\plotplaceholder}[2]{
    \begin{figure}[H]
        \centering
        \IfFileExists{../outputs/lab3/plots/#1}{
            \includegraphics[width=0.88\textwidth]{#1}
        }{
            \fbox{
                \parbox[c][5cm][c]{0.82\textwidth}{
                    \centering
                    Файл рисунка \texttt{\detokenize{#1}} создается запуском\\
                    \texttt{python lab3/run.py}
                }
            }
        }
        \caption{#2}
    \end{figure}
}

\title{Лабораторная работа №3\\Модификации градиентного спуска}
\author{
    Новиков Алексей Георгиевич, группа М3233\\
    Душкин Александр Алексеевич, группа М3233
}
\date{2026}

\begin{document}

\begin{titlepage}
    \begin{center}
        {\large Отчет по лабораторной работе №3}\\[0.5cm]
        {\LARGE \textbf{Модификации градиентного спуска}}\\[2cm]

        \begin{flushright}
            Выполнили:\\[0.2cm]
            Новиков Алексей Георгиевич, группа М3233\\
            Душкин Александр Алексеевич, группа М3233\\[0.5cm]
        \end{flushright}

        \vfill
        Санкт-Петербург\\[0.2cm]
        2026
    \end{center}
\end{titlepage}

\tableofcontents
\newpage

\section{Постановка задачи}

Цель работы --- реализовать и исследовать модификации градиентного спуска, которые используют инерцию, накопление статистики градиентов или адаптивное масштабирование шага по координатам.

В работе реализованы:
\begin{enumerate}
    \item Momentum;
    \item Nesterov accelerated gradient;
    \item AdaGrad;
    \item RMSProp;
    \item AdaDelta;
    \item Adam.
\end{enumerate}

Во всех экспериментах использовалась фиксированная точность
\[
    \varepsilon = 10^{-8},
\]
а критерием остановки была норма градиента:
\[
    \|\nabla f(x_k)\|\le \varepsilon.
\]

\textbf{Промежуточный вывод.} Третья лабораторная продолжает вторую: базовое движение по антиградиенту заменяется правилами, которые используют память о прошлых шагах или статистику прошлых градиентов. Поэтому основной объект исследования смещается к гиперпараметрам: именно они определяют, насколько сильно метод доверяет текущему градиенту, прошлому направлению и накопленным масштабам по координатам.

\section{Используемые модели задач}

Модели задач переиспользуются из лабораторной работы №2:
\begin{itemize}
    \item хорошо обусловленная квадратичная функция
    \[
        f(x,y)=\frac{1}{2}(x^2+y^2);
    \]
    \item плохо обусловленная квадратичная функция
    \[
        g(x,y)=\frac{1}{2}x^2+50y^2;
    \]
    \item функция Розенброка
    \[
        r(x,y)=(1-x)^2+100(y-x^2)^2;
    \]
    \item функция Экли;
    \item функция Химмельблау.
\end{itemize}

Для квадратичных функций использовалась стартовая точка $(-2,-2)$. Для сложных функций в конфигурации задавались несколько стартовых точек, чтобы проверить чувствительность методов к области притяжения.

\textbf{Промежуточный вывод.} Повторное использование моделей из второй лабораторной делает сравнение корректным: меняется не задача, а только правило построения следующего приближения. Это важно методически, потому что Momentum, RMSProp и Adam отличаются не направлением базового градиента как таковым, а способом преобразовать его в обновление с учетом истории. На одних и тех же поверхностях можно увидеть, какие трудности решает память метода, а какие остаются из-за геометрии функции.

\section{Описание реализованных методов}

\subsection{Momentum}

Метод Momentum добавляет инерционный член:
\[
    v_t=\beta v_{t-1}+\alpha\nabla f(x_t),
\]
\[
    x_{t+1}=x_t-v_t.
\]
Параметр $\beta$ отвечает за сохранение предыдущего направления движения.

\textbf{Промежуточный вывод.} Инерция ускоряет движение вдоль устойчивого направления, потому что новый шаг складывается не только из текущего градиента, но и из накопленной скорости. Если несколько градиентов подряд имеют близкое направление, скорость растет и метод быстрее проходит пологие участки. Но та же логика может навредить около минимума или в узкой долине: накопленная скорость продолжает тянуть точку вперед, даже когда текущий градиент уже требует торможения или поворота.

\subsection{Метод Нестерова}

В методе Нестерова градиент вычисляется в точке предпросмотра:
\[
    \tilde{x}_t=x_t-\beta v_{t-1},
\]
\[
    v_t=\beta v_{t-1}+\alpha\nabla f(\tilde{x}_t),
\]
\[
    x_{t+1}=x_t-v_t.
\]

\textbf{Промежуточный вывод.} Метод Нестерова пытается заранее оценить, куда приведет инерция, поэтому градиент вычисляется не в текущей точке, а в точке предпросмотра. Такая структура делает корректировку направления более ранней: метод видит последствия накопленной скорости до фактического обновления. За это он платит дополнительным вычислением градиента на каждой итерации, что видно в счетчиках вызовов.

\subsection{AdaGrad}

AdaGrad накапливает сумму квадратов градиентов:
\[
    s_t=s_{t-1}+\nabla f(x_t)^2,
\]
\[
    x_{t+1}=x_t-\alpha\frac{\nabla f(x_t)}{\sqrt{s_t}+\varepsilon_{\mathrm{num}}}.
\]

\textbf{Промежуточный вывод.} AdaGrad уменьшает эффективный шаг по координатам, где градиенты часто велики, потому что знаменатель содержит накопленную сумму квадратов градиента. Это полезно при разных масштабах координат: направление автоматически нормируется по истории. Однако накопитель только растет, поэтому на длинных детерминированных запусках эффективный шаг монотонно уменьшается и метод может почти остановиться до достижения требуемой точности.

\subsection{RMSProp}

RMSProp заменяет полную сумму квадратов экспоненциальным средним:
\[
    s_t=\rho s_{t-1}+(1-\rho)\nabla f(x_t)^2,
\]
\[
    x_{t+1}=x_t-\alpha\frac{\nabla f(x_t)}{\sqrt{s_t}+\varepsilon_{\mathrm{num}}}.
\]

\textbf{Промежуточный вывод.} Экспоненциальное сглаживание предотвращает бесконечный рост накопителя, потому что старые квадраты градиента постепенно забываются с коэффициентом $\rho$. В отличие от AdaGrad, RMSProp масштабирует шаг по недавней, а не по всей накопленной истории. Поэтому метод лучше адаптируется, когда масштаб градиента меняется по мере движения к минимуму.

\subsection{AdaDelta}

AdaDelta использует экспоненциальные средние квадратов градиентов и квадратов обновлений:
\[
    E[g^2]_t=\rho E[g^2]_{t-1}+(1-\rho)g_t^2,
\]
\[
    \Delta x_t=-\frac{\sqrt{E[\Delta x^2]_{t-1}+\varepsilon_{\mathrm{num}}}}
    {\sqrt{E[g^2]_t+\varepsilon_{\mathrm{num}}}}g_t,
\]
\[
    E[\Delta x^2]_t=\rho E[\Delta x^2]_{t-1}+(1-\rho)\Delta x_t^2.
\]

\textbf{Промежуточный вывод.} Метод не требует явного глобального шага, потому что масштаб обновления определяется отношением накопленного среднего квадратов прошлых шагов к среднему квадратов градиентов. Такая самонормировка делает AdaDelta устойчивым к выбору learning rate, но в рассматриваемых задачах приводит к консервативным обновлениям: если средний размер прошлых шагов мал, новые шаги также остаются малыми, и метод часто доходит до лимита итераций.

\subsection{Adam}

Adam объединяет Momentum и RMSProp:
\[
    m_t=\beta_1m_{t-1}+(1-\beta_1)g_t,
\]
\[
    v_t=\beta_2v_{t-1}+(1-\beta_2)g_t^2.
\]
Используется коррекция смещения:
\[
    \hat{m}_t=\frac{m_t}{1-\beta_1^t}, \qquad
    \hat{v}_t=\frac{v_t}{1-\beta_2^t}.
\]
Обновление:
\[
    x_{t+1}=x_t-\alpha\frac{\hat{m}_t}{\sqrt{\hat{v}_t}+\varepsilon_{\mathrm{num}}}.
\]

\textbf{Промежуточный вывод.} Adam является наиболее универсальным из рассмотренных адаптивных методов, потому что объединяет две разные идеи: первый момент $m_t$ задает сглаженное направление движения, а второй момент $v_t$ нормирует шаг по координатам. Коррекция смещения особенно важна в начале оптимизации, когда накопители еще малы и без нее первые обновления были бы систематически искажены.

\section{Структура программной реализации}

Лабораторная работа построена поверх той же библиотеки \texttt{optimlib}, что и лабораторная №2. Ключевые элементы:
\begin{itemize}
    \item \texttt{AdaptiveOptimizer} --- базовый класс для методов с внутренним состоянием;
    \item \texttt{\_reset\_state()} --- инициализация буферов перед каждым запуском;
    \item \texttt{\_adaptive\_direction(...)} --- метод, в котором конкретный оптимизатор строит вектор обновления;
    \item \texttt{Momentum}, \texttt{Nesterov}, \texttt{AdaGrad}, \texttt{RMSProp}, \texttt{AdaDelta}, \texttt{Adam} --- реализации методов;
    \item \texttt{OptimizationExperiment} --- перебор функций, стартовых точек и сеток параметров;
    \item \texttt{plot\_param\_sensitivity} --- построение line-графиков и heatmap чувствительности.
\end{itemize}

Сетки гиперпараметров задаются в \texttt{lab3/config.yaml}. Для Momentum и Nesterov перебираются $\alpha$ и $\beta$, для RMSProp --- $\alpha$ и $\rho$, для AdaGrad --- $\alpha$, для AdaDelta --- $\rho$, для Adam --- $\alpha$, $\beta_1$ и $\beta_2$.

\textbf{Промежуточный вывод.} Общая архитектура позволила добавить шесть новых оптимизаторов без изменения кода функций и экспериментального раннера, потому что различия методов локализованы в построении вектора обновления. Базовый класс отвечает за цикл оптимизации, критерии остановки и историю, а конкретные методы хранят только свои состояния: скорость, накопитель квадратов градиента или моменты Adam. Такая структура хорошо соответствует математике методов с памятью.

\section{Результаты вычислительных экспериментов}

\subsection{Квадратичные функции}

В таблице приведены лучшие наборы параметров по числу итераций на квадратичных функциях.

\begin{table}[H]
\centering
\caption{Лучшие параметры на квадратичных функциях при $\varepsilon=10^{-8}$}
\begin{tabular}{llrrrrr}
\toprule
Функция & Метод & $\alpha$ & $\beta/\rho$ & $\beta_1$ & $\beta_2$ & Итерации\\
\midrule
$f$ & Momentum & $10^{-2}$ & $0.85$ & -- & -- & 182\\
$f$ & Nesterov & $10^{-2}$ & $0.85$ & -- & -- & 202\\
$f$ & AdaGrad & $10^{-3}$ & -- & -- & -- & 1500\\
$f$ & RMSProp & $10^{-2}$ & $0.85$ & -- & -- & 238\\
$f$ & AdaDelta & -- & $0.50$ & -- & -- & 1500\\
$f$ & Adam & $10^{-2}$ & -- & $0.70$ & $0.90$ & 289\\
$g$ & Momentum & $10^{-2}$ & $0.85$ & -- & -- & 260\\
$g$ & Nesterov & $10^{-2}$ & $0.85$ & -- & -- & 202\\
$g$ & AdaGrad & $10^{-3}$ & -- & -- & -- & 1500\\
$g$ & RMSProp & $10^{-2}$ & $0.85$ & -- & -- & 239\\
$g$ & AdaDelta & -- & $0.50$ & -- & -- & 1500\\
$g$ & Adam & $10^{-2}$ & -- & $0.70$ & $0.95$ & 362\\
\bottomrule
\end{tabular}
\end{table}

\textbf{Промежуточный вывод.} На квадратичных функциях наиболее быстрыми оказались Momentum, Nesterov и RMSProp. У инерционных методов это объясняется накоплением направления: на квадратичной поверхности градиенты меняются достаточно предсказуемо, поэтому скорость помогает проходить пологие участки. RMSProp выигрывает за счет координатного масштабирования, полезного особенно на плохо обусловленной функции. AdaGrad и AdaDelta в выбранной сетке часто не достигали точности до лимита, потому что их механизмы уменьшения эффективного шага делали обновления слишком осторожными.

\subsection{Чувствительность к параметрам}

Для методов с двумя параметрами строились heatmap, где цвет отражает число итераций. Для методов с одним параметром строились line-графики.

\plotplaceholder{WellConditionedQuadratic__x0_-2.0_-2.0__Momentum_heatmap.png}{Чувствительность Momentum к параметрам $\alpha$ и $\beta$}
\plotplaceholder{WellConditionedQuadratic__x0_-2.0_-2.0__RMSProp_heatmap.png}{Чувствительность RMSProp к параметрам $\alpha$ и $\rho$}
\plotplaceholder{WellConditionedQuadratic__x0_-2.0_-2.0__Adam_heatmap.png}{Чувствительность Adam к параметрам $\beta_1$ и $\beta_2$}
\plotplaceholder{WellConditionedQuadratic__x0_-2.0_-2.0__AdaGrad_sensitivity.png}{Чувствительность AdaGrad к начальному шагу}

\textbf{Промежуточный вывод.} Для инерционных методов важна согласованность шага и коэффициента инерции: $\alpha$ задает силу текущего градиента, а $\beta$ определяет, сколько прошлой скорости переносится в новый шаг. Если оба эффекта велики, обновление может стать слишком агрессивным и начать колебаться. Для адаптивных методов аналогичную роль играет сглаживание: малое сглаживание быстро реагирует на текущий градиент, а большое делает масштаб шага стабильнее, но медленнее приспосабливается к новой области поверхности.

\subsection{Траектории на квадратичных функциях}

\plotplaceholder{WellConditionedQuadratic__x0_-2.0_-2.0__best_trajectories.png}{Лучшие траектории методов на хорошо обусловленной квадратичной функции}
\plotplaceholder{IllConditionedQuadratic__x0_-2.0_-2.0__best_trajectories.png}{Лучшие траектории методов на плохо обусловленной квадратичной функции}

\textbf{Промежуточный вывод.} На хорошо обусловленной функции траектории в основном направлены к центру линий уровня, потому что масштаб по координатам одинаков и направление градиента хорошо согласовано с направлением к минимуму. На плохо обусловленной функции адаптивное масштабирование и инерция уменьшают влияние вытянутости разными способами: инерция накапливает движение вдоль устойчивого направления, а RMSProp и Adam уменьшают шаг по координатам с большими недавними градиентами. Поэтому они выигрывают у простого градиентного спуска с малым шагом.

\subsection{Сложные функции}

В таблице приведены лучшие результаты по значению функции среди исследованных параметров и стартовых точек.

\begin{table}[H]
\centering
\caption{Лучшие запуски на сложных функциях}
\begin{tabular}{lllr}
\toprule
Функция & Метод & Лучший набор параметров & Итерации\\
\midrule
Ackley & Momentum & $\alpha=3\cdot10^{-4},\ \beta=0.95$ & 562\\
Ackley & Nesterov & $\alpha=3\cdot10^{-3},\ \beta=0.50$ & 43\\
Ackley & AdaGrad & $\alpha=10^{-2}$ & 81\\
Ackley & RMSProp & $\alpha=10^{-4},\ \rho=0.90$ & 376\\
Ackley & AdaDelta & $\rho=0.85$ & 175\\
Ackley & Adam & $\alpha=3\cdot10^{-4},\ \beta_1=0.70,\ \beta_2=0.95$ & 230\\
Rosenbrock & Momentum & $\alpha=3\cdot10^{-3},\ \beta=0.90$ & 1277\\
Rosenbrock & Nesterov & $\alpha=10^{-3},\ \beta=0.95$ & 1500\\
Rosenbrock & Adam & $\alpha=10^{-2},\ \beta_1=0.85,\ \beta_2=0.99$ & 1500\\
Himmelblau & Momentum & $\alpha=3\cdot10^{-3},\ \beta=0.50$ & 63\\
Himmelblau & Nesterov & $\alpha=10^{-2},\ \beta=0.50$ & 24\\
Himmelblau & Adam & $\alpha=10^{-3},\ \beta_1=0.85,\ \beta_2=0.99$ & 1339\\
\bottomrule
\end{tabular}
\end{table}

\textbf{Промежуточный вывод.} На функции Химмельблау методы обычно находили один из глобальных минимумов, потому что после попадания в область притяжения градиентные направления достаточно уверенно ведут к ближайшему минимуму. На функции Розенброка ситуация сложнее: узкая изогнутая долина требует не только правильного масштаба шага, но и постоянного изменения направления вдоль кривой. Инерция может ускорять движение, но также может уносить поперек долины; адаптивные масштабы помогают с длиной шага, но не устраняют полностью проблему направления.

\subsection{Траектории на сложных функциях}

\plotplaceholder{Rosenbrock__x0_-1.2_1.0__best_trajectories.png}{Лучшие траектории методов на функции Розенброка}
\plotplaceholder{Ackley__x0_1.0_1.0__best_trajectories.png}{Лучшие траектории методов на функции Экли}
\plotplaceholder{Himmelblau__x0_2.0_2.0__best_trajectories.png}{Лучшие траектории методов на функции Химмельблау}

\textbf{Промежуточный вывод.} Визуализация траекторий особенно важна для сложных функций: одинаковое итоговое число итераций может соответствовать совершенно разному поведению. У инерционных методов на графиках видны более длинные и сглаженные перемещения, поскольку шаг зависит от прошлой скорости. У адаптивных методов траектория может быть менее прямой, зато координатное масштабирование уменьшает слишком резкие движения в направлениях с большими градиентами. Поэтому график траектории дополняет таблицу и показывает, как именно внутренняя логика метода влияет на путь к минимуму.

\section{Сводная таблица эксперимента}

Полная сводка всех $1470$ запусков сохраняется в файле \texttt{../outputs/lab3/summary.csv}. В ней содержатся функция, стартовая точка, оптимизатор, параметры запуска, признак сходимости, итоговая точка, значение функции и счетчики вычислений.

\IfFileExists{../outputs/lab3/summary.tex}{
    \scriptsize
    \input{../outputs/lab3/summary.tex}
    \normalsize
}{
    \begin{center}
        \fbox{\parbox{0.82\textwidth}{Полная таблица доступна в \texttt{../outputs/lab3/summary.csv}.}}
    \end{center}
}

\section{Общий вывод}

В лабораторной работе были реализованы шесть модификаций градиентного спуска. Эксперименты показали, что добавление инерции и адаптивного масштабирования может существенно ускорить сходимость по сравнению с классическим градиентным спуском, но делает методы чувствительными к гиперпараметрам.

Momentum и Nesterov хорошо проявили себя на квадратичных функциях. Для обеих квадратичных моделей удачными оказались параметры порядка $\alpha=10^{-2}$ и $\beta=0.85$. Метод Нестерова иногда требовал меньше итераций на плохо обусловленной функции, однако платил за это дополнительными вычислениями градиента.

RMSProp показал устойчивую работу на квадратичных функциях и подтвердил полезность экспоненциального сглаживания квадратов градиента. AdaGrad и AdaDelta в выбранной сетке часто двигались слишком осторожно и не достигали точности до лимита итераций.

Adam оказался универсальным методом, особенно полезным как компромисс между инерцией и адаптивным шагом. При этом на детерминированных тестовых задачах он не всегда был самым быстрым: на простых квадратичных функциях Momentum и RMSProp часто требовали меньше итераций.

На сложных функциях результат сильнее зависел от стартовой точки и геометрии поверхности. Функция Розенброка осталась наиболее трудной из-за узкой изогнутой долины, а функция Химмельблау хорошо решалась из стартовых точек, попадающих в области притяжения глобальных минимумов.

\end{document}
